%



%%%

%%%   Самарский государственный технический университет

%%%   Кафедра прикладной математики и информатики

%%%

%%%   Преамбула для статьи в журнал "Вестник Самарского государственного

%%%   технического университета. Серия: «Физико-математические науки»"

%%%   Ориентирована под MikTEx 2.4 for Win32

%%%

%%%   Хотя сам журнал собирается под GNU/Linux на дистрибутиве TeXLive



\documentclass[11pt,a4paper,twoside]{article}



\usepackage[cp1251]{inputenc}

\usepackage[T2A]{fontenc}

\usepackage[english,russian]{babel}

\usepackage{samgtubib}

\usepackage[dvips]{graphicx}

\usepackage[multiple]{footmisc}



\usepackage{samgtu}

\renewcommand{\VestnikYear}{2009} % Год издания Вестника

\renewcommand{\VestnikNumber}{2\,(19)} % Номер Вестника

\renewcommand{\VestnikMonth}{Ноябрь} % Месяц выхода Вестника

\renewcommand{\VestnikData}{01.11.09} % Дата подписания Вестника в печать

\renewcommand{\VestnikUPL}{26,64} % Усл. печ. л.

\renewcommand{\VestnikUIL}{25,73} % Уч.-изд. л.

\renewcommand{\VestnikRegNumber}{479/09} % Регистрационный номер

\renewcommand{\VestnikOrder}{994} % Номер заказа









%%%%

%%%   Самарский государственный технический университет

%%%   Кафедра прикладной математики и информатики

%%%

%%%   Преамбула для статьи в журнал "Вестник Самарского государственного

%%%   технического университета. Серия: «Физико-математические науки»"

%%%   Ориентирована под MikTEx 2.4 for Win32

%%%

%%%   Хотя сам журнал собирается под GNU/Linux на дистрибутиве TeXLive



\documentclass[11pt,a4paper,twoside]{article}



\usepackage[cp1251]{inputenc}

\usepackage[T2A]{fontenc}

\usepackage[english,russian]{babel}

\usepackage{samgtubib}

\usepackage[dvips]{graphicx}

\usepackage[multiple]{footmisc}



\usepackage{samgtu}

\renewcommand{\VestnikYear}{2009} % Год издания Вестника

\renewcommand{\VestnikNumber}{2\,(19)} % Номер Вестника

\renewcommand{\VestnikMonth}{Ноябрь} % Месяц выхода Вестника

\renewcommand{\VestnikData}{01.11.09} % Дата подписания Вестника в печать

\renewcommand{\VestnikUPL}{26,64} % Усл. печ. л.

\renewcommand{\VestnikUIL}{25,73} % Уч.-изд. л.

\renewcommand{\VestnikRegNumber}{479/09} % Регистрационный номер

\renewcommand{\VestnikOrder}{994} % Номер заказа









\usepackage{qrcode}



%

%\newcommand{\totop}[1]{\raisebox{6pt}[1pt][2pt]{#1}} 

%\newcommand{\tottop}[1]{\raisebox{12pt}[1pt][2pt]{#1}} 

%\newcommand{\totttop}[1]{\raisebox{18pt}[1pt][2pt]{#1}} 

%\newcommand{\tottttop}[1]{\raisebox{24pt}[1pt][2pt]{#1}} 

%\newcommand{\totttttop}[1]{\raisebox{30pt}[1pt][2pt]{#1}} 

\newcommand{\tobot}[1]{\raisebox{-6pt}[1pt][2pt]{#1}} 

%\newcommand{\tobbot}[1]{\raisebox{-12pt}[1pt][2pt]{#1}} 

%\newcommand{\tobbbot}[1]{\raisebox{-18pt}[1pt][2pt]{#1}}



%\graphicspath{{./00PIC/}}



%\samgtupubl % для генерации pdf, отдаваемого в типографию СамГТУ

\pdfcrop % обрезка pdf для размещения на math-net.ru

\draft % На корректуре

\filllastpage % Последняя страница не заполнена не полностью

%\briefcomm % Статья является кратким сообщением



\vsgtu{1508}



\FirstPage{1}

\LastPage{24}

%

%

\begin{document}





\renewcommand{\baselinestretch}{0.9}



\hfill \qrcode[hyperlink,height=.8in]{http://dx.doi.org/10.14498/vsgtu1508}

\vspace{-.8in}





% \Partition{Механика деформируемого твердого тела}{Applied mechanics} % Раздел Вестника, в который подаётся статья

% Названия разделов см. на сайте при подаче заявки





%%%%%%%%%%%%%%%%%%%%%%%%%%%%%%%%%%%%%%%%%%%%%%%%%%%%%%%%%%%%%%%%%%%%%%%%%%%%%%%%%%%%

%Информация о статье и об авторах на русском и английском языках



% Номер УДК

\udc{???}%Обработка текста. Подготовка текстов : Языки программирования. Компьютерные языки

\msc{???}% Коды MSC см. http://www.ams.org/msc/



\rutitle{Математическое моделирование напряженно-деформированного состояния размеростабильных композитных элементов конструкций оптических телескопов с помощью метода конечных элементов}% Полный заголовок

{Математическое моделирование напряженно-деформированного состояния \dots}% Заголовок, который идёт в верхний колонтитул



\ruauthor{В.\,Е.~Биткин\inst1, О.\,Г.~Жидкова\inst1, А.\,В.~Денисов\inst1, А.\,В.~Бородавин\inst1, Д.\,В.~Митюшкина\inst1, А.\,В.~Родионов\inst1, А.\,С.~Нонин\inst2}% Автор(ы) для заголовка, если авторов несколько укать \inst

{Б\,и\,т\,к\,и\,н~В.\,Е.,\ Ж\,и\,д\,к\,о\,в\,а~О.\,Г.,\ Д\,е\,н\,и\,с\,о\,в~А.\,В.,\ Б\,о\,р\,о\,д\,а\,в\,и\,н~А.\,В.,\ М\,и\,т\,ю\,ш\,к\,и\,н\,а~Д.\,В.,\ Р\,о\,д\,и\,о\,н\,о\,в~А.\,В.,\ Н\,о\,н\,и\,н~А.\,С.}% Автор(ы) для оглавления и колонтитула через \,



\ruaffil{ООО <<СКТБ <<Пластик>>,\\Россия, 446025, Самарская обл., г. Сызрань, Саратовское шоссе, 4. 

\and 

АО РКЦ <<Прогресс>>,\\Россия, 443009, Самара, ул.Земеца, 18.}

\enaffil{LLC SKTB Plastic,\\4, Saratov hwy., Syzran, 446025, Russian Federation. 

\and 

JSC SRC Progress,\\18, Zemetsa str., Samara, 443009, Russian Federation.}

% Если несколько, то так  \institute{ Организация 1 \and Организация 2  \and Организация 3}



%\email{E-mails}{petuhovds@mail.ru, kie@icmm.ru} %E-mail's авторов



\ruauthordetails{7}{\fullauthor{\rupid{123028}{Владимир Евгеньевич Биткин}} \regalia{\mem{gksi@sktb-plastik.ru}; автор, ведущий переписку} \position{первый заместитель генерального директора "--- генеральный конструктор по СИ}\\

\fullauthor{\rupid{123029}{Ольга Геннадьевна Жидкова}} \regalia{\mem{opriokr-prg@sktb-plastik.ru}}\position{заместитель генерального конструктора по научной

работе}\\

\fullauthor{\rupid{123030}{Александр Владимирович Денисов}} \regalia{\mem{opriokr@sktb-plastik.ru}}\position{начальник отдела}\\

\fullauthor{\rupid{123031}{Андрей Викторович Бородавин}} \regalia{\mem{opriokr-zn3@sktb-plastik.ru}}\position{заместитель начальника отдела}\\

\fullauthor{\rupid{123032}{Диана Викторовна Митюшкина}} \regalia{\mem{opriokr-prg@sktb-plastik.ru}}\position{ведущий инженер}\\

\fullauthor{\rupid{123033}{Александр Вениаминович Родионов}} \regalia{\mem{grp1@sktb-plastik.ru}}\position{руководитель проектов}\\

\fullauthor{\rupid{123034}{Александр Сергеевич Нонин}} \regalia{\mem{1104@samspace.ru}}\position{заместитель начальника отдела "--- начальник сектора}  }



\enauthordetails{7}{\fullauthor{\enpid{123028}{Vladimir E. Bitkin}} \regalia{\mem{gksi@sktb-plastik.ru}; Corresponding Author} \position{First Deputy General Director --- General Designer for measuring systems}\\

\fullauthor{\enpid{123029}{Olga G. Zhidkova}} \regalia{\mem{opriokr-prg@sktb-plastik.ru}} \position{Deputy General Designer for Scientific Research}\\

\fullauthor{\enpid{123030}{Alexander V. Denisov}} \regalia{\mem{opriokr@sktb-plastik.ru}} \position{Head of Department}\\

\fullauthor{\enpid{123031}{Andrey V. Borodavin}} \regalia{\mem{opriokr-zn3@sktb-plastik.ru}} \position{Deputy Head of Department}\\

\fullauthor{\enpid{123032}{Diana V. Mityushkina}} \regalia{\mem{opriokr-prg@sktb-plastik.ru}} \position{Senior Engineer}\\

\fullauthor{\enpid{123033}{Alexander V. Rodionov}} \regalia{\mem{grpl@sktb-plastik.ru}} \position{Project Manager}\\

\fullauthor{\enpid{123034}{Alexander S. Nonin}} \regalia{\mem{1104@samspace.ru}} \position{Deputy Head of Department --- Division Head}}





\rukeywords{крупногабаритные композитные конструкции, принцип интегральности композитных конструкций, размеростабильная несущая конструкция, оптический телескоп, параметрический анализ конструкции, метод конечных элементов, математическое моделирование, этапы проектирования размеростабильных конструкций}



\ruabstract{Рассмотрены вопросы проектирования размеростабильных композитных элементов конструкций оптических телескопов. Описана этапность проведения проектировочных расчётов. Приведены основные соотношения микромеханики теории композитных материалов. На примере разработанной композитной конструкции корпуса оптико"=электронного комплекса с односторонним подкреплением ребрами описаны особенности математического моделирования с учётом принятых допущений. Представлены результаты экспериментального определения характеристик углепластиков, используемые при проектировании размеростабильной несущей конструкции корпуса оптического телескопа, отражены преимущества метода конечных элементов как одного из основных методов решения краевых задач прикладной механики. Показана корректность использования аналитических методов на начальных этапах разработки с целью сокращения сроков проектирования. Определена ведущая роль конечно"=элементного моделирования в прогнозировании поведения конструкций на различных этапах эксплуатации. Показаны размеростабильные несущие композитные корпуса, разработанные с учётом описанной последовательности проектирования конструкций. Описанная этапность создания конструкции позволила обрабатывать и систематизировать данные в процессе их разработки и экспериментальной отработки, уточнять параметры модели конструкции при параметрическом анализе.}



\entitle{Mathematical simulation for strain-stress stste of optical telescope stable-size composite elements with finite-element method}

%

\enauthor{V.~Е.~Bitkin\inst1, О.~G.~Zhidkova\inst1, А.~V.~Denisov\inst1, А.~V.~Borodavin\inst1, D.~V.~Mityushkina\inst1, А.~V.~Rodionov\inst1, A.~S.~Nonin\inst2}{B\,i\,t\,k\,i\,n~V.~Е., Z\,h\,i\,d\,k\,o\,v\,a~О.~G., D\,e\,n\,i\,s\,o\,v~А.~V., B\,o\,r\,o\,d\,a\,v\,i\,n~А.~V., M\,i\,t\,y\,u\,s\,h\,k\,i\,n\,a~D.~V., R\,o\,d\,i\,o\,n\,o\,v~А.~V., N\,o\,n\,i\,n~A.~S.}





\enabstract{Designing issues of optical telescope stable-size composite elements have been considered. Staging of design calculation for composite structural elements has been described. Basic relations of composite material micro-mechanics have been represented. Specifics of mathematical simulation taking into account assumptions made have been described by example of developed framework of electro-optical system with one-side enforcement by ribs. Results of experimental determination of carbon-filled plastic characteristics used in design of stable-size optical telescope frame structure have been represented; advantages of finite-element method as one of the basic methods for solving the boundary value problem in applied mechanics have been reflected. Reasonableness of analytical approach using in the initial development stage in order to shorten the period of design has been demonstrated. The leading part of finite-element simulation has been determined in behavior prognostication of structures at different operating stages. Stable-size supporting composite frameworks developed taking into account defined sequence of structural design have been showed. Described staging of structure making has been allowed to process and systematize data during design and experimental execution, refine structural model parameters, increase the confidence level and verify it. }



\enkeywords{large-size composite structures, composite structure integrity concept, stable-size supporting structure, optical telescope, parametric analysis of structure, finite-element method, mathematical simulation, staging of stable-size structures}



\date{13/VII/2016}

\revisiondate{26/VIII/2016}

\accepteddate{09/IX/2016}



%%%%%%%%%%%%%%%%%%%%%%%%%%%%%%%%%%%%%%%%%%%%%%%%%%%%%%%%%%%%%%%%%%%%%%%%%%%%%%%%%%%%





\makerutitle





\Section[n]{Введение} При проектировании крупногабаритных композитных космических конструкций к ним предъявляется целый ряд требований, которые обусловлены наличием механических нагрузок, термических напряжений, ограничениями допустимого изменения размеров и необходимостью получения минимальной массы всей конструкции \cite{bit:bib1}. Эффективная работа таких крупногабаритных конструкций, как конструкции космических оптических телескопов, во многом определяется правильным выбором концепции построения и минимизацией влияния искажающих факторов на функционирование их систем. Концепция построения телескопа в первую очередь зависит как от выбора оптической схемы, обеспечивающей формирование оптического излучения и изображения с минимальными потерями энергии и информации, так и от выбора конструкции, отвечающей требованиям по прочностным, жесткостным и стабилизационным характеристикам в реальных условиях действия телескопа.



Основная сложность проблемы разработки таких конструкций заключается в том, что эффективное формообразование космических конструкций телескопов может быть определено как синтез формостабилизированных систем с заданным законом деформирования. При этом формостабилизация обеспечивается в динамике процесса эксплуатационных изменений напряжённо"=деформированного состояния конструкций \cite{bit:bib2}.



В качестве искажающих факторов в конструкциях космических телескопов выступают дифракция, фазировка элементов, аберрации оптической системы, термодеформация оправ и конструкции, тепловые неоднородности в пространстве, стабильность юстировки оптической системы, деформация элементов и конструкции, ошибки изготовления, системы контроля и др.



Анализ погрешностей показывает, что не только правильность выбора оптической схемы, её технологичность и коррекционные возможности вносят основной вклад в качество изображения, но и стабильность размеров несущих конструкций телескопов и оправ зеркал играет существенную роль в получении качественных изображений.



При создании таких сложных космических конструкций для более глубокого понимания сущности явлений деформирования должны быть с наибольшей достоверностью выполнены натурные эксперименты. Однако, испытание таких конструкций в условиях Земли может оказаться физически невозможным. Поэтому одной из неотложных задач становится разработка надёжных математических методов прогнозирования поведения такого рода конструкций в космических условиях. Усилия специалистов в настоящее время сосредоточены на оптимизации массы, разработке метода моделирования поведения крупногабаритных конструкций космического назначения в основном методом конечных элементов и изучения нелинейных характеристик больших конструкций \cite{bit:bib3}.



Принимая во внимание надёжность имеющихся ЭВМ, точность расчёта характеристик конструкций определяется качеством физических и математических моделей. Из-за сложности конструкций определение всех их характеристик часто не может основываться на результатах экспериментальных испытаний, особенно, если эти конструкции являются крупногабаритными \cite{bit:bib4, bit:bib5} (рис.~\ref{bit:fig1}). При использовании армированных слоистых композиционных материалов проблема усложняется наличием присущих им остаточных напряжений, появляющихся в процессе отверждения или в результате действия механических нагрузок и неблагоприятных внешних условий (влажности, температуры).



\begin{figure}[h!]

\centering

\footnotesize

\includegraphics[width=0.47\textwidth]{image1-1}

\hfill

\includegraphics[width=0.47\textwidth]{image1-2}

\\

{\sl а} \hspace{70mm} {\sl б}

\caption{Современные конструкции крупногабаритных космических телескопов: {\sl а} "--- <<Хаббл>>; {\sl б} "--- <<Кеплер>>\label{bit:fig1}}[Figure~\ref{bit:fig1}. ???]

\smallskip  

\end{figure}



Надёжность расчётов повышается, если предварительно найдена корреляция между теоретическими и экспериментальными данными для отдельных узлов. Убедительное доказательство надёжности теоретических расчётов может быть получено только после их сравнения с результатами экспериментальных испытаний. Для крупногабаритных космических конструкций подобная проверка практически невозможна из-за очень больших размеров и стоимости таких конструкций. Тем не менее, основные задачи программы испытаний могут быть решены на сборных узлах. Вопрос об анализе поведения крупногабаритных конструкций, которые невозможно испытать на Земле из-за их размеров или необходимости создания условий невесомости, всё ещё остается открытым. В настоящее время в связи с интенсивным развитием вычислительной техники и численных методов расширились возможности исследований статической и динамической прочности сложных инженерных конструкций, таких как, например, космические конструкции оптико"=электронных комплексов. Однако не всегда можно достигнуть удовлетворительного решения, полагаясь лишь на параметры высокопроизводительных вычислительных комплексов. Основу процессов исследования и проектирования конструкций составляют все-таки модели деформирования, от которых в значительной степени зависит достижение качеств, настойчиво предъявляемых практикой проектирования, таких, как повышение достоверности, информативности и оперативности исследований \cite{bit:bib6}.



Возможности развития такого рода космических конструкций нового класса неразрывно связаны с процессом в области создания новых композиционных материалов. Стремление получить наивысшие эксплуатационные характеристики привело к разработке и использованию в космических конструкциях композиционных материалов "--- высокомодульных углепластиков нового поколения на основе цианат"=эфирного связующего с минимальным содержанием летучих конденсирующихся веществ (ЛКВ), а также минимальными значениями пористости, влагопоглощения и общей потери массы.



При совмещении процессов формования элементов конструкции со сборкой в неотверждённом или полуотверждённом состояниях и последующим совместным отверждением, необходимость в выполнении механических соединений почти полностью отпадает. В результате сокращаются затраты времени на сборку и уменьшается суммарная себестоимость. Такой подход к проектированию, изготовлению и сборке изделий привел к разработке принципа интегральности композитных конструкций, позволяющего уменьшить число входящих элементов (за счёт большой жёсткости и прочности композиционного материала) \cite{bit:bib7}.



Интегральные конструкции из композиционных материалов "--- новый прогрессивный вид технологии (в широком смысле), исходящий из учёта специфических технологических и механических свойств композитов и позволяющий создавать конструкции не только с меньшей массой, но и с меньшей себестоимостью изготовления сборки, чем металлические аналоги \cite{bit:bib8}.



\Section[n]{Этапы проектирования и роль конечно"=элементного моделирования в создании конструкций оптических телескопов} Разрабатываемые конструктивно"=технологические решения по созданию высокоточных интегральных конструкций из КМ являются сложной и фундаментальной задачей, затрагивающей не только методы изготовления, но и свойства исходных материалов, расчёт и проектирование конструкций, условия эксплуатации и учёта дополнительных требований, накладывающих ограничения по точности и чистоте рабочей поверхности \cite{bit:bib2}.



Сложность решаемых задач требует постоянного совершенствования процесса проектирования конструкций на базе современного математического обеспечения с учётом опыта использования отечественных и зарубежных систем автоматизированного проектирования (САПР) или CAD (Computer-Aided Design) и систем автоматизации инженерных расчётов и анализа (CAE), позволяющих моделировать, в частности, процессы деформации, исчерпания несущей способности и разрушения конструкций сложной конфигурации, работающих в экстремальных условиях эксплуатации.



Одним из основополагающих звеньев при проектировании конкретного вида конструкции оптического телескопа является анализ напряжённо"=деформированного состояния с выбором конструкционных материалов, определением основных расчётных случаев нагружения, отвечающий за весовое совершенство проектируемого изделия \cite{bit:bib6}.



\begin{figure}[h!]

\centering

\footnotesize

\includegraphics[width=0.7\textwidth]{image2}

\caption{Элементы интегральной размеростабильной композитной конструкции несущего каркаса космического оптико"=электронного комплекса: 1)~проставка коническая; 2)~корпус объектива\label{bit:fig2}}[Figure~\ref{bit:fig2}. ???]

\smallskip  

\end{figure}



Необходимость специализированного математического обеспечения возникает при использовании в конструкции космических аппаратов композиционных материалов, в частности, углепластика, которые создаются в процессе изготовления самой конструкции. Так, проектирование опытной размеростабильной несущей конструкции оптического телескопа (рис.~\ref{bit:fig2}) из полимерных композиционных материалов сопровождалось исследованием напряжённо"=деформированного состояния его конструктивных элементов и осуществлялось с помощью программного комплекса конечно"=элементного анализа MSC.Nastran \cite{bit:bib9}.



Именно для конструкций подобного типа надёжное математическое обеспечение, особенно предиспытательное математическое моделирование, имеет высокий статус в связи с возможностью значительного облегчения корпусных элементов оптического телескопа, а также уменьшением (в ряде случаев) объёма натурных повторных испытаний при условии обеспечения углубленного деформативного и прочностного анализа.



В настоящее время обеспечение размеростабильности космических конструкций реализует комплексный подход, состоящий из двух частей: расчётно"=теоретической и экспериментальной, и базирующейся на двух основных принципах. Первый "--- это нормирование уровня температурных и силовых нагрузок, второй "--- проведение испытаний на специальных экспериментальных стендах.



Следует отметить, что проведение прямого расчёта на прочность и деформативность конструкции при температурном и силовом воздействии "--- задача сложная и трудоёмкая даже для разработанной конструкции, облик которой определён чертёжной документацией или 3D-моделью. В качестве проектного расчёта при разработке конструкции в такой постановке задача не только нецелесообразна из-за неопределённости исходных параметров, но и порой невыполнима \cite{bit:bib3}.



В целом процесс по обеспечению прочности и деформативности на статические температурные и силовые нагрузки можно разделить на следующие этапы.



\underline {Первый этап.} По упрощённой модели, исходя из паспортных данных на материал наполнителя и связующего, проводится термомеханический анализ свойств монослоя композиционного материала. Из полученных термомеханических характеристик однонаправленного материала конструируется слоистый композиционный материал, определяются через обобщённые жёсткости композита его средние модули упругости, коэффициенты термического линейного расширения, коэффициенты Пуассона, коэффициенты влияния в зависимости от рассматриваемых схем армирования монослоёв. Далее путём проектировочных расчётов на силовые и температурные нагрузки проводится разработка конструкции с назначением размеров, уточняются схемы армирования и свойства композита, определяющие прочность и жёсткость конструкции. В заключение проводится проверочный расчёт.



Следует отметить, что наиболее эффективными в отношении удельных характеристик, являются композиты, образованные из непрерывных волокон и полимерной матрицы. Основная задача для исследователя заключается в вычислении эффективных модулей упругости, которые определяются как коэффициенты, связывающие усреднённые напряжения и деформации (рис.~\ref{bit:fig3}).



\begin{figure}[h!]

\centering

\footnotesize

\includegraphics[width=0.6\textwidth]{bit-fig3}

\caption{Элемент однонаправленного композита\label{bit:fig3}}[Figure~\ref{bit:fig3}. ???]

\smallskip  

\end{figure}



Для ориентировочных оценок и качественного анализа целесообразно использовать соотношения влияния микроструктурных параметров на свойства композита в соответствии с законом механической смеси \cite{bit:bib10}. Тогда для одного слоя (см. рис. 3) имеем следующие соотношения \cite{bit:bib2}:

\begin{gather*}

%\label{bit:eq1}

%\begin{array}{c}

E_1=E_\textit{в}\nu+E_\textit{м}(1-\nu);

\\

\displaystyle \frac{1}{E_2}=\frac{\nu}{E_\textit{в}}+\frac{1-\nu}{E_\textit{м}}+\frac{[\mu_\textit{в}\nu+\mu_\textit{м}(1-\nu)]^2}{E_\textit{в}\nu+E_\textit{м}(1-\nu)};

\\

\displaystyle \frac{1}{G_{12}}=\frac{\nu}{G_\textit{в}}+\frac{1-\nu}{G_\textit{м}};

\\

\mu_{21}=\mu_\textit{в}\nu+\mu_\textit{м}(1-\nu);

\\

E_1\mu_{12}=E_2\mu_{21};

\\

\displaystyle \alpha_1=\frac{E_\textit{в}\alpha_\textit{в}\nu+E_\textit{м}\alpha_\textit{м}(1-\nu)}{E_\textit{в}\nu+E_\textit{м}(1-\nu)};

\\

\displaystyle \alpha_2=\alpha_\textit{в}(1+\mu_\textit{в})\nu+\alpha_\textit{м}(1+\mu_\textit{м})(1-\nu)-

\\

-[\mu_\textit{в}\nu+\mu_\textit{м}(1-\nu)]\frac{E_\textit{в}\alpha_\textit{в}\nu+E_\textit{м}\alpha_\textit{м}(1-\nu)}{E_\textit{в}\nu+E_\textit{м}(1-\nu)},

%\end{array}

\end{gather*}

где $E_1$ "--- продольный модуль упругости одного слоя; $E_{2}$ "--- поперечный модуль упругости слоя; $G_{12}$ "--- средний модуль сдвига слоя; $\alpha_1$ "--- коэффициент линейного термического расширения (КЛТР) материала вдоль волокон; $\alpha_2$ "--- КЛТР материала поперёк волокон; $\mu_{12}$, $\mu_{21}$ "--- коэффициенты Пуассона слоя. $E_{\textit{м}}$, $E_{\textit{в}}$, G$_{\textit{м}}$, $G_{\textit{в}}$, $\alpha_{\textit{м}}$, $\alpha _{\textit{в}}$ "--- модули упругости, модули сдвига, коэффициенты линейного термического расширения (КЛТР) матрицы и волокон.



Соотношения, позволяющие выразить упругие постоянные слоистого композита (пакета из армированных слоёв, см. рис.~\ref{bit:fig4}) через его структурные параметры, имеют вид \cite{bit:bib2, bit:bib11}:

\begin{gather*}

%\label{bit:eq2}

%\begin{array}{c}

\displaystyle E_x=\frac{Q}{Q_{22}Q_{33}-Q_{23}^2};

\\

\displaystyle E_y=\frac{Q}{Q_{11}Q_{33}-Q_{13}^2};

\\

\displaystyle G_{xy}=\frac{Q}{Q_{11}Q_{22}-Q_{12}^2};

\\

\displaystyle \mu_{xy}=\frac{Q_{12}Q_{33}-Q_{13}Q_{23}}{Q_{22}Q_{33}-Q_{23}^2};

\\

\displaystyle \mu_{yx}=\frac{Q_{12}Q_{33}-Q_{13}Q_{23}}{Q_{11}Q_{33}-Q_{13}^2};

\\

\displaystyle \eta_{x,xy}=\eta_{xy,x}=\frac{Q_{12}Q_{23}-Q_{22}Q_{13}}{Q_{11}Q_{22}-Q_{12}^2};

\\

\displaystyle \eta_{y,xy}=\eta_{xy,y}=\frac{Q_{12}Q_{13}-Q_{11}Q_{23}}{Q_{11}Q_{22}-Q_{13}^2};

\\

\displaystyle \alpha_x=\frac{1}{E_xh}(Q_{1T}-\mu_{yx}Q_{2T}+\eta_{x,xy}Q_{3T});

\\

\displaystyle \alpha_y=\frac{1}{E_yh}(Q_{2T}-\mu_{xy}Q_{1T}+\eta_{y,xy}Q_{3T});

\\

\displaystyle \alpha_{xy}=\frac{1}{G_{xy}h}(\eta_{xy,x}Q_{1T}+\eta_{xy,y}Q_{2T}+Q_{3T}),

%\end{array}

\end{gather*}

где $E_x$, $E_y$, $G_{xy}$ "--- осреднённые модули упругости, модуль сдвига слоистого КМ; $\mu_{xy}$, $\mu_{yx}$ "--- осреднённые коэффициенты Пуассона слоистого КМ; $\eta_{x,xy}$, $\eta_{xy,x}$, $\eta_{y,xy}$, $\eta_{xy,y}$ "--- осреднённые коэффициенты влияния II рода слоистого пакета КМ; $\alpha_x$, $\alpha_y$, $\alpha_{xy}$ "--- осреднённые КЛТР слоистого пакета КМ; $Q_{ij}$, $Q_{iT}$ "--- обобщённые жёсткости слоистого КМ ($i,j=1,2,3$), которые имеют вид:

\begin{gather*}

Q_{11}=\sum_{i=0}^{k}h_i(E_1^i\cos^4\theta_i+2E_1^i\mu_{12}^i\sin^2\theta_i\cos^2\theta_i+E_2^i\sin^4\theta_i+G_{12}^i\sin^22\theta_i);

\\

Q_{21}=Q_{12}=\sum_{i=0}^{k}[(E_1^i+E_2^i)\sin^2\theta_i\cos^2\theta_i+E_1^i\mu_{12}^i(\sin^4\theta_i+\cos^4\theta_i)-G_{12}^i\sin^22\theta_i];

\\

Q_{22}=\sum_{i=0}^{k}h_i(E_1^i\sin^4\theta_i+2E_1^i\mu_{21}^i\sin^2\theta_i\cos^2\theta_i+E_2^i\sin^4\theta_i+G_{12}^i\sin^22\theta_i);

\\

Q_{31}=Q_{13}=\sum_{i=0}^{k}h_i\sin\theta_i\cos\theta_i[E_1^i(1-\mu_{21}^i)\cos^2\theta_i-E_2^i(1-\mu_{12}^i)\sin^2\theta_i-2G_{12}^i\cos2\theta_i];

\\

Q_{23}=Q_{32}=\sum_{i=0}^{k}h_i\sin\theta_i\cos\theta_i[E_1^i(1-\mu_{21}^i)\sin^2\theta_i-E_2^i(1-\mu_{12}^i)\cos^2\theta_i-2G_{12}^i\cos2\theta_i];

\\

Q_{33}=\sum_{i=0}^{k}[(E_1^i+E_2^i-2E_1^i\mu_{21}^i)\sin^2\theta_i\cos^2\theta_i+G_{21}^i\cos^22\theta_i];

\\

Q_{1T}=\sum_{i=0}^{k}h_i[E_1^i(\alpha_1^i+\mu_{21}^i\alpha_2^i)\cos^2\theta_i+E_2^i(\alpha_2^i+\mu_{21}^i\alpha_1^i)\sin^2\theta_i];

\\

Q_{2T}=\sum_{i=0}^{k}h_i[E_1^i(\alpha_1^i+\mu_{21}^i\alpha_2^i)\sin^2\theta_i+E_2^i(\alpha_2^i+\mu_{12}^i\alpha_1^i)\cos^2\theta_i];

\\

Q_{3T}=\sum_{i=0}^{k}h_i\sin\theta_i\cos\theta_i[E_1^i(\alpha_1^i+\mu_{21}^i\alpha_2^i)-E_2^i(\alpha_2^i+\mu_{12}^i\alpha_1^i)],

\end{gather*}

где $\theta_i$ "--- угол армирования монослоя композита.



\begin{figure}[h!]

\centering

\footnotesize

\includegraphics[width=0.6\textwidth]{bit-fig4}

\caption{Элемент слоистого композиционного материала\label{bit:fig4}}[Figure~\ref{bit:fig4}. ???]

\smallskip  

\end{figure}



На рисунке \ref{bit:fig5} показаны расчётные графические зависимости интегральных характеристик композиционного материала в зависимости от угла армирования $\theta$ (см. рис.~\ref{bit:fig4}) для углепластика на основе углеродной ленты ЛУ-П/0.1 и связующего ЭНФБ  со степенью армирования (объёмное содержание углеродных волокон в пакете) $V_{\textit{в}}=0.6$ и схемой армирования $(0^\circ/\pm\theta^\circ/90^\circ)_{2n}$. Индекс $n$ (натуральное число) показывает, что число слоёв в пакете чётное. Отметим, что толщина всего пакета $h$ может достигать величины до 5~мм, а толщина одного слоя пакета $h_1$ составляет $0.1-0.13$~мм (см. рис.~\ref{bit:fig4}). Анализ данных зависимостей показывает, что, варьируя углы армирования слоёв пакета $\theta$, т.~е. используя чередующуюся схему укладки слоёв с углами $\pm\theta$, можно получить термомеханические характеристики композита, позволяющие свести к минимуму как деформации конструкции от воздействия температуры, так и создать двуоснотермонейтральные структуры пакета КМ. Эти возможности могут быть в значительной степени расширены путём использования методов технологического воздействия, выбора режима отверждения, степени армирования композита, технологического напряжения и др.



\begin{figure}[h!]

\vspace{-3mm}

\centering

\footnotesize

\includegraphics[width=0.8\textwidth]{bit-fig5}

\caption{Зависимость КЛТР и модуля упругости от угла армирования композита $\theta$ для углепластика на основе углеродной ленты ЛУ-П/0.1 и связующего ЭНФБ

(степень армирования $V_{\textit{в}}=0.6$, схема армирования $(0^\circ/\pm\theta^\circ/90^\circ)_{2n}$)\label{bit:fig5}}[Figure~\ref{bit:fig5}. ???]

\smallskip  

\end{figure}



При проектировании конструкции космического оптического телескопа и его элементов, таких, как размеростабильные несущие конструкции корпуса, идёт процесс синтеза формы будущей конструкции и материала, из которого она будет изготовлена \cite{bit:bib3, bit:bib12}. При этом форма конструкции имеет большое значение для обеспечения необходимой её жёсткости и прочности. В связи с этим возникает необходимость в надёжных методах расчёта конструкций любой формы.



В настоящее время конечно"=элементное моделирование поведения конструкции в условиях эксплуатации является неотъемлемой частью процесса проектирования инженерно и структурно сложных объектов, к которым относятся изделия из композиционных материалов. Актуальность и сложность задачи построения адекватных математических моделей многократно возрастают при создании высокоточных размеростабильных композитных конструкций с крайне малыми допусками по деформационным параметрам.



\underline {Второй этап.} Для уточнения расчётных значений термомеханических характеристик многослойного композиционного материала "--- углепластика, проводятся экспериментальные исследования, включающие изготовление плоских образцов композита с однонаправленными слоями и целесообразно выбранными укладками, определяются коэффициенты вариации модуля упругости, предела прочности, коэффициента линейного термического расширения и др., проводится оценка весовой эффективности, определяются поправочные коэффициенты. Следует отметить, что полученные экспериментальным путём характеристики слоистого композита позволяют провести анализ напряжённо"=деформированного состояния конструкции из реальных материалов \cite{bit:bib13}. Такой подход позволяет более точно спрогнозировать поведение конструкции в условиях силового и температурного воздействий. При проектировании корпуса размеростабильной несущей конструкции оптико"=электронного комплекса были проведены исследования термомеханических характеристик углепластиков в исходном состоянии при температуре $20^\circ С$. Использовались три типа углеродного наполнителя "--- ленты ЛУ-П/0.1 (ГОСТ 28006-88), Кулон-500/0.07 (СТО 75969490-007-2009) и углеродное волокно М55J (Япония). В качестве связующих использовались эпоксиноволачноформальдегидное ЭНФБ (ТУ 1-596-36-2005) и цианат"=эфирное НИИКАМ-ЦЭМОС (ТУ 2242-060-14527989-2012). В таблице 1 представлены результаты для однонаправленных углепластиков и углепластиков со схемой армирования \mbox{$(0^\circ/\pm45^\circ/90^\circ)_{2n}$}. Коэффициенты вариации модуля упругости $K_{Е}$ и предела прочности $K_{\sigma }$ определялись в процентах соотношением:

$$K=\frac{\sqrt{\frac{1}{n}\sum_{i=1}^{n}(x_{\textit{ср}}-x_i)^2}}{x_{\textit{ср}}}\cdot100,$$

где $x_i$ "--- экспериментальное значение исследуемой величины для $i$-го образца; $x_{\textit{ср}}$ "--- среднее значение этой же величины.



\begin{table}[h!]

\centering\footnotesize

\caption{Экспериментальные значения характеристик углепластиков [???]\label{bit:tabl1}}

\vspace{-3mm}

\begin{tabular}[h]{|p{8mm}|p{8mm}|p{6mm}|p{6mm}|p{5mm}|p{8mm}|p{6mm}|p{6mm}|p{8mm}|p{4mm}|p{5mm}|p{20mm}|}%{|p{0.15\linewidth}|p{0.20\linewidth}|p{0.2\linewidth}| p{0.20\linewidth}|}%{|c|c|c|c||c|c|}

 \hline 

%\footnotesize Номер& \tobot{\footnotesize $f$, Гц} & \footnotesize Номер & \tobot{\footnotesize $f$, Гц} & \footnotesize Номер & \tobot{\footnotesize $f$, Гц} \\

%\footnotesize частоты &  & \footnotesize частоты &  & \footnotesize частоты &  \\

\tiny Материал & \tiny Укладка & \tiny Направление вырезки & \tiny Модуль упругости при растяжении $E_{\textit{р}}$, МПа & \tiny Предел прочности при растяжении $\sigma _{\textit{р}}$, МПа & \tiny Коэффициент Пуассона ($\nu$) & \tiny Модуль упругости при сжатии $E_{\textit{сж}}$, МПа & \tiny Предел прочности при изгибе $\sigma_{\textit{и}}$, (МПа) & \tiny Модуль упругости при изгибе $E_{\textit{и}}$, МПа & \tiny Предел прочности на сдвиг $\tau$, МПа & \tiny Модуль сдвига $G_{12}$, МПа & \tiny Коэффициент линейного термического расширения $\alpha\cdot10^{-6}$, $K^{-1}$ при $20^\circ$C \\

\hline

\raisebox{-4.50ex}[0cm][0cm]{\tiny КМУ-4Л \par (ЛУ-П/0.1 на связующем ЭНФБ)}& \raisebox{-1.50ex}[0cm][0cm]{\tiny однонаправленная}& \tiny 0\r{ }& \tiny 137255& \tiny 637& \tiny 0.251& -& \tiny 1078& \tiny 153922& -& -& \tiny $-0.39\div-0.53$ \\

\cline{3-12}

 & & \tiny $90^\circ$&-&-&-&-&-&-&-&-& \tiny $32.3\div35.6$ \\

\cline{2-12}

 & \raisebox{-1.50ex}[0cm][0cm]{\tiny $(0^\circ$/$\pm45^\circ$/$90^\circ)_{2n}$}& \tiny $0^\circ$ & \tiny 56785& \tiny 270& \tiny $0.33$ & -& \tiny $196$& -& -& -& \tiny $1.26\div1.41$ \\

\cline{3-12}

 &  & \tiny $90^\circ$ & \tiny 56785& \tiny 285& \tiny $0.33$ & -& -& -& \tiny 135& \tiny 24501& \tiny $1.25\div1.43$ \\

\hline

\raisebox{-1.50ex}[0cm][0cm]{\tiny Кулон-500/0.07 на связующем ЭНФБ}&

\raisebox{-1.50ex}[0cm][0cm]{\tiny однонаправленная}&

\tiny $0^\circ$ & \tiny 245166& -& \tiny $0.302$ & -& -& -& -& \tiny 4903& \tiny $-0.5\div-0.7$ \\

\cline{3-12}

 &  & \tiny $90^\circ$ & \tiny 5884& -& \tiny $0.00703$ & -& -& -& -& -& \tiny 34 \\

\hline

\raisebox{-4.50ex}[0cm][0cm]{\tiny М55J/ \par НИИКАМ-РС}&

\raisebox{-1.50ex}[0cm][0cm]{\tiny однонаправленная}&

\tiny $0^\circ$ & \tiny 298910& \tiny 1395& \tiny $0.35$& \tiny 401383& \tiny 851& \tiny 201077& -& -& \tiny $-0.864\div-0.891$ \\

\cline{3-12}

 &  & \tiny $90^\circ$ & -& -& -& -& -& -& -& -& \tiny $29.6\div29.8$ \\

\cline{2-12}

 & \raisebox{-1.50ex}[0cm][0cm]{\tiny $(0^\circ$/$\pm45^\circ$/$90^\circ)_{2n}$}&

\tiny $0^\circ$& \tiny 121409& \tiny 537& \tiny $0.31$& \tiny 172086& \tiny 561& \tiny 91984& -& -& \tiny $-0.138\div-0.098$ \\

\cline{3-12}

 &  & \tiny $90^\circ$& \tiny 107144& \tiny 586& \tiny $0.33$ & -& -& -& \tiny 172& \tiny 34870& \tiny $0.026\div-0.030$ \\

\hline

\raisebox{-4.50ex}[0cm][0cm]{\tiny ЛУ-П/0.1 и цианат"=эфирное связующее НИИКАМ-ЦЭМОС}&

\raisebox{-1.50ex}[0cm][0cm]{\tiny однонаправленная}&

\tiny $0^\circ$& \tiny 160227& \tiny 896& \tiny $0.27$& \tiny 216067& \tiny 1155& \tiny 114518& -& -& \tiny $-0.300\div-0.206$ \\

\cline{3-12}

 &  & \tiny $90^\circ$& -& -& -& -& -& -& -& -& \tiny $24.0\div23.5$ \\

\cline{2-12}

 &

\raisebox{-1.50ex}[0cm][0cm]{\tiny $(0^\circ$/$\pm45^\circ$/$90^\circ)_{2n}$}&

\tiny $0^\circ$& \tiny 65298& \tiny 268& \tiny $0.32$& \tiny 74257& \tiny 485& \tiny 55063& -& -& \tiny $1.24\div1.36$ \\

\cline{3-12}

 &  & \tiny $90^\circ$& \tiny 71851& \tiny 384& \tiny $0.33$& -& -& -& \tiny 148& \tiny 25583& \tiny $1.22\div1.35$ \\

\hline

\end{tabular}

\end{table}



Полученные экспериментальные значения модуля упругости, модуля сдвига хорошо коррелируют с расчётными. Экспериментальные данные по всем параметрам имеют незначительный разброс. В частности, значения коэффициента вариации находятся в пределах от 4.05 до 14.6{\%}. Отклонения же расчётных значений модулей упругости, коэффициентов Пуассона, коэффициентов линейного термического расширения и других параметров от соответствующих экспериментальных значений не превосходят 10\%.



\underline {Третий этап.} На основе полученных на втором этапе параметров строится уточнённая конечно"=элементная модель (КЭМ) конструкции, адекватно отражающая её жёсткостные и прочностные характеристики, достаточно подробная для проведения напряжённо"=деформированного анализа. Эта модель является инструментом для проведения повторного анализа с целью получения достоверного поведения конструкции с учётом элементов навесного оборудования. В рамках этой модели должно быть получено НДС конструкции, соответствующее нормативным силовым и температурным нагрузкам, действующим в условиях эксплуатации. На основе последних можно определить предельное состояние конструкции и, соответственно, запасы по деформативности и прочности.



Теоретический анализ напряжённо"=деформированного состояния позволяет создать эталон деформированной поверхности размеростабильной несущей конструкции оптических телескопов в условиях нагружения тепловыми и механическими нагрузками \cite{bit:bib14}. Численное моделирование поведения несущих композитных корпусов телескопов в условиях орбитальной эксплуатации и при наземной юстировке показывает, что при использовании существующих в настоящее время материалов с их реальными термомеханическими характеристиками удаётся достичь некоторого порогового минимума отклонений угловых и линейных размеров за счёт подбора рациональной структуры армирования несущих слоёв. Однако с увеличением габаритных размеров размеростабильного корпуса возрастает и пороговое значение минимальных отклонений.



При решении задач по определению напряжённо"=деформированного состояния с использованием численного анализа на базе метода конечных элементов для достижения необходимой степени точности важное значение имеет выбор моделирующего элемента. Кроме ошибок, вызванных вычислительной процедурой, метод конечных элементов имеет ошибки с построением математической модели исследуемого объекта. К ним относятся ошибки дискретизации, ошибки, связанные различием реальной границы моделируемого объекта и её аппроксимации, ошибки выбора базисной функции, обусловленной разностью между точным решением и его представлением некоторой функцией в пределах элемента \cite{bit:bib15, bit:bib16}.



В связи с тем, что аналитические критерии сходимости той или иной задачи МКЭ весьма сложны и представляют собой область отдельных исследований, в качестве критерия сходимости численного решения к точному решению может быть использована экспериментальная проверка результатов расчёта.



\underline {Четвёртый этап.} Проведённый анализ и разработанная модель позволяют подготовить экспериментальную отработку несущей способности изготовленной конструкции, определить максимальные температурные деформации и выделить параметры исследуемой модели, подлежащие проверке и уточнению \cite{bit:bib1}.



В процессе экспериментальной отработки проверяются и подтверждаются результаты анализа напряжённо"=деформированного состояния конструкции, проводятся изменения, уточняющие и подтверждающие основные характеристики конструкции.



\underline {Пятый этап.} Проводится анализ результатов измерений, полученных в процессе испытаний (рис.~\ref{bit:fig6}), корректировка параметров конструкции. При необходимости проводятся расчёты, моделирующие процесс нагружения конструкции при испытаниях, подтверждается степень достоверности откорректированной модели.



\begin{figure}[h!]

\vspace{-3mm}

\centering

\footnotesize

\reflectbox{\includegraphics[width=0.95\textwidth]{image6}}

\caption{Испытание изделия в термо"=вакуумной камере\label{bit:fig6}}[Figure~\ref{bit:fig6}. ???]

\smallskip  

\end{figure}



\underline {Шестой этап.} Подтверждённая результатами испытаний конечно"=элементная модель может быть использована при изменённых условиях эксплуатации для достоверных расчётов изделия, внесения изменений в конструкцию нового изделия подобного класса.



Основная роль на всех этапах работ в рамках действующей системы отработки термодеформаций и прочности принадлежит модели. При этом развитие системы неизбежно пойдет по двум направлениям \cite{bit:bib6}:



-- переход нагрузок к действительным, соответствующим уровням эксплуатации;



-- перенос экспериментальных работ со специализированного макета на лётное изделие.



В процессе прогнозируемого развития роль модели будет возрастать, как и требования к ней по полноте, подробности, достоверности.



В настоящее время метод конечных элементов в силу своих многочисленных достоинств вытеснил практически все другие методы решения краевых задач прикладной механики; при использовании данного метода, в отличие от классических методов рассматриваются свойства элементов конечных размеров, а интегрирование по объёму конструкции заменяется конечным суммированием интегралов по объёмам элементов. Уравнения в частных производных для сплошной среды сводятся к обыкновенным дифференциальным или алгебраическим уравнениям \cite{bit:bib15}.



Конечно"=элементная модель является математическим отображением конструкции и её параметров: геометрических характеристик, упругих свойств конструкционных материалов и др. Опыт проведения инженерных расчётных и экспериментальных работ по созданию размеростабильных композитных корпусов оптических телескопов позволяет выделить основные направления разработки конечно"=элементной модели изделия оптико"=электронного комплекса:



-- конечно"=элементная модель описывается линейными стержневыми и двумерными пластинчатыми и оболочечными элементами \cite{bit:bib17}; использование трёхмерных объёмных элементов, приводящих к значительному увеличению размерности КЭМ, должно быть обосновано;



-- моделирование стыков элементов и закрепления дискретных вводимых в модель навесных элементов конструкции осуществляется с помощью интегральных и специальных соединительных элементов \cite{bit:bib18};



-- степень дискретности разбиения на конечные элементы по размерам выбирается исходя из необходимости обеспечения адекватного отражения деформативных и прочностных характеристик до определённого значения; размеры элементов в зонах высокоградиентного напряжённого состояния должны выявлять места, опасные с точки зрения разрушения \cite{bit:bib19};



-- геометрические, упругие характеристики вводятся в модель номинальными значениями, изменения или модификации модели должны быть легко реализуемы.



Основные трудности, которые приходится преодолевать разработчику конечно"=элементной модели, связаны со специфическими особенностями нетрадиционных конструкторских решений, контактными задачами на стыке элементов конструкции телескопа, введением условных конечных элементов, отображающих интегральные характеристики совокупности соединённых между собой или взаимодействующих элементов конструкции. Особого внимания и подхода требуют геометрически изменяемые конструкции и их элементы.



В настоящее время для проведения анализа напряжённо"=деформированного состояния и постановки экспериментальных работ используются универсальные профессиональные пакеты прикладных программ конечноэлементного анализа ANSYS, Nastran и др. \cite{bit:bib17, bit:bib20, bit:bib21}.



Современный уровень программных средств позволяет в полной мере реализовать указанные выше принципы и преодолеть сложности, обеспечивая возможность использования набора стандартных конечных элементов. В процессе расчёта осуществляются операции с матрицами высокой размерности ленточной структуры, размерность матриц определяется количеством узлов и типом применяемых конечных элементов. При этом полный цикл задач термоупругости предполагает многократное проведение различных видов ресурсоемких расчётов. Кроме того, в последнее время хорошие результаты по производительности демонстрируются при распараллеливании вычислений с использованием современных процессоров.



Однако, для детального анализа локального НДС в критических местах с точки зрения деформативности и прочности, приходится использовать дополнительные, более локальные модели, определяя из общей модели граничные условия и силовые факторы, необходимые для расчётов локальной модели.



Конечноэлементная модель является основным инструментом расчётного решения по теоретическому определению запасов прочности конструкции при силовых и температурных воздействиях. Результатом анализа является вычисление напряжённо"=деформированного состояния конструкции, в соответствии со спецификой которого определяется её предельное состояние.



Математическое моделирование определяет наиболее нагруженные места конструкции размеростабильного несущего корпуса, кронштейнов, других элементов крепления. Это позволяет организовать измерение деформаций, установив в определённых зонах конструкции тензодатчики. Конечно"=элементная модель играет основную роль при разработке изделия в решении задач прочности, жёсткости, температурных деформаций на всех этапах создания конструкции, в том числе и на стадии экспериментальной отработки в процессе испытаний. Предметом измерения, уточнения и экспериментального подтверждения являются важнейшие характеристики конструкции: жёсткостные, прочностные, весовые и пр. Результатом анализа несоответствий между расчётными и экспериментальными значениями характеристик конструкции являются мероприятия по корректировке модели, направленные на сокращение величины рассогласований.



Инженерная практика показывает, что хорошим соответствием можно считать расхождение в пределах 10{\%} по деформативным характеристикам. Точность результатов практически полностью определяется соответствием интегральных характеристик материала условным характеристикам конечных элементов реальной конструкции. Настройка конечно"=элементной модели по деформативной картине производится варьированием жёсткостных параметров конечных элементов. В некоторой степени настройка может быть уточнена за счёт изменения геометрических параметров элементов конструкции, заданных первоначально номинальными значениями. Сложность в подборе жёсткостных характеристик заключается в том, что изменение параметров одного элемента модели влечёт изменение других частей модели, в связи с этим настройка конечно"=элементной модели требует итерационного подхода и связана с необходимостью проведения многократных расчётов, моделирующих деформативную картину поведения конструкции.



\begin{figure}[h!]

\vspace{-3mm}

\centering

\footnotesize

\includegraphics[width=0.325\textwidth]{image7-1}\hfill\includegraphics[width=0.41\textwidth]{image7-2}\hfill\includegraphics[width=0.242\textwidth]{image7-3}

\caption{Размеростабильные элементы несущих конструкций корпусов оптико"=электронных комплексов космических аппаратов\label{bit:fig7}}[Figure~\ref{bit:fig7}. ???]

\smallskip  

\end{figure}



Рассмотренная последовательность создания размеростабильной конструкции была использована при разработке размеростабильных несущих композитных корпусов телескопов (рис.~\ref{bit:fig7}), позволяла обрабатывать и систематизировать данные в процессе разработки и экспериментальной отработки, настраивать модель, уточнять её параметры, поднять уровень её достоверности, адекватно отражать внутренние свойства конструкции и её реакцию на внешние воздействия \cite{bit:bib22}.



\Section[n]{Результаты численного моделирования полей деформаций размеростабильных элементов конструкций космических телескопов} При анализе и конечно"=элементном моделировании деформативного поведения углепластикового размеростабильного несущего корпуса объектива оптико"=электронного комплекса, представляющего собой интегральную композитную конструкцию с односторонним подкреплением с внешней стороны продольно"=поперечным набором рёбер, были сделаны следующие допущения \cite{bit:bib15}:



-- для оболочки справедливы гипотезы Кирхгофа"--~Лява;



-- для рёбер применима теория стержней Кирхгофа"--~Клебша с учётом продольной деформации, изгиба и кручения;



-- зависимость между напряжениями и деформациями отвечают закону Гука;



-- модули упругости, коэффициенты Пуассона и коэффициенты термического линейного расширения "--- величины постоянные;



-- температура по толщине обшивки и высоте рёбер может меняться по линейному закону.



При исследовании геометрии деформированной конструкции предпочтение отдано функционалу Лагранжа, который построен на базе перемещений $U$, $V$, $W$ с учётом принятых гипотез и имеет форму:

$$

\text{Э}(U,V,W)=\frac12\sum_{i=1}^{n}\sum_{j=1}^{k}\text{Э}^\text{об}(U,V,W)+\frac12\sum_{i=1}^{n}\sum_{j=1}^{k}\text{Э}^\text{р}(U,V,W).

$$



Функционал представлен в виде суммы энергий всех гладких прямоугольных областей ($\text{Э}^\text{об}$) и рёбер ($\text{Э}^\text{р}$), окаймляющих их.



Результатами параметрического анализа являются выбор свойств композиционного материала, количества рёбер продольно"=поперечного набора, определение массы, перемещений и углов поворота конструкции, а также напряжённого состояния в процессе эксплуатации.



Конечно"=элементное моделирование поведения конструкции размеростабильного несущего корпуса из углепластика для оптико"=электронного комплекса было проведено с использованием программного комплекса конечно"=элементного анализа. Оболочка корпуса моделировалась оболочечными элементами, учитывающими все внутренние силовые факторы. Рёбра жёсткости аппроксимировались балочными элементами, прикрепляемыми к узлам разбивки с эксцентриситетом. Для моделирования влияния навешиваемого оборудования (элементов оптической системы) использовались элементы массы, моделирующие поступательную и вращательную инерцию, сосредоточенную в узле. Для передачи нагрузки использовались элементы жёсткой связи и балочные элементы. Крепление конструкции моделировалось с учётом материала ответной части, также изготовленной из углепластика с целью исключения пространственного деформирования при тепловом нагружении. Корпус закреплён по стыковочному фланцу по всем поступательным степеням свободы, за исключением радиального направления, в цилиндрической системе координат, ось $OZ$ которой совпадает с продольной осью конструкции. Однако для численной реализации методом конечных элементов использовалась декартова система координат. На рисунке~\ref{bit:fig8} приведена конечно"=элементная модель несущей конструкции корпуса оптико"=электронного комплекса (ОЭК) в декартовой системе координат, начало которой находится в центре нижнего крепёжного фланца, а ось $OZ$ совпадает с продольной осью корпуса.



\begin{figure}[h!]

\vspace{-3mm}

\centering

\footnotesize

\includegraphics[width=0.5\textwidth]{image8}

\caption{Конечно"=элементная модель несущей конструкции корпуса оптико"=электронного комплекса\label{bit:fig8}}[Figure~\ref{bit:fig8}. ???]

\smallskip  

\end{figure}



В соответствии с требованиями, предъявляемыми к конструкции корпуса по размерной стабильности при воздействии теплового нагружения при штатной эксплуатации, проведён анализ характера деформирования корпуса при изменении температуры $T_{\text{ном}}\pm5{ }^\circ$С, где $T_{\text{ном}}=20{ }^\circ$С. При этом вследствие идентичности поведения конструкции при <<отрицательных>> и <<положительных>> перепадах температур, оценивались линейные и угловые отклонения элементов оптической системы при $\Delta T=T_{\text{ном}}+5{ }^\circ$С. Для наихудшего расчётного случая с точки зрения деформативности конструкции (неравномерный прогрев в диаметральном направлении в плоскости $XOY$) построены расчётные зависимости линейных смещений $\Delta X$, $\Delta Y$, $\Delta Z$ и разворотов вокруг оси $OY$ элементов оптической системы (рис.~\ref{bit:fig9}--\ref{bit:fig10}) для конструкции корпуса из углепластика на основе углелент ЛУ-П/0.1 и Кулон-500/0.07 на связующем ЭНФБ, схема армирования $(0^\circ/\pm45^\circ/90^\circ)_{2n}$. Угловые смещения относительно осей $OX$, $OZ$ имеют малые значения ($\ll1$ угл. сек.) и практически не оказывают влияния на картину деформативности корпуса при воздействии указанной выше нагрузки. На рис.~\ref{bit:fig9} и~\ref{bit:fig10} величина $\Delta T$ "--- это разность температур на внешней и внутренней границах оболочки в диаметральной плоскости $XOY$. В расчётах принималось, что температура по толщине изменялась по линейному закону при заданной разности $\Delta T$.



\begin{figure}[h!]

\vspace{-3mm}

\centering

\footnotesize

\includegraphics[width=0.47\textwidth]{bit-fig9-1}\hfill\includegraphics[width=0.47\textwidth]{bit-fig9-2}

\caption{Зависимость смещений $\Delta X$, $\Delta Y$, $\Delta Z$ элементов оптической системы от перепада температур для конструкции корпуса из углепластика на основе углелент ЛУ-П/0.1 и Кулон-500/0.07\label{bit:fig9}}[Figure~\ref{bit:fig9}. ???]

\smallskip  

\end{figure}



\begin{figure}[h!]

\vspace{-3mm}

\centering

\footnotesize

\includegraphics[width=0.47\textwidth]{bit-fig10}

\caption{Зависимость разворотов относительно оси $ОY$ элементов оптической системы от перепада температур для конструкции корпуса из углепластика на основе углелент ЛУ-П/0.1 и Кулон-500/0.07\label{bit:fig10}}[Figure~\ref{bit:fig10}. ???]

\smallskip  

\end{figure}



Картина суммарных перемещений корпуса (углепластик на основе углеленты ЛУ-П/0.1 на связующем ЭНФБ) при воздействии перепада температур в диаметральном направлении $\Delta T=T_{\text{ном}}+5^\circ$С, где $T_{\text{ном}}=20^\circ$С, показаны на рисунке~\ref{bit:fig11}.



\begin{figure}[h!]

\centering

\footnotesize

\includegraphics[width=0.47\textwidth]{image11}

\caption{Картина суммарных перемещений конструкции корпуса из углепластика на основе углеленты ЛУ-П/0.1 и связующего ЭНФБ при воздействии перепада температур в диаметральном направлении $\Delta T=T_{\text{ном}}+5^\circ$С ($T_{\text{ном}}=20^\circ$С).\label{bit:fig11}}[Figure~\ref{bit:fig11}. ???]

\smallskip  

\end{figure}



Анализ графических зависимостей, представленных на рисунках~\ref{bit:fig9},~\ref{bit:fig10}, показывает незначительное преимущество с точки зрения деформативности при температурном нагружении для конструкции корпуса из материала на основе углеленты Кулон-500/0.07. При этом требуемые значения стабильности пространственных линейных и угловых положений посадочных площадок под элементы оптической системы в условиях штатной эксплуатации оптико"=электронного комплекса достигаются и для конструкции корпуса из углепластика на основе углеленты ЛУ-П/0.1. Необходимо отметить, что углепластик на основе углеленты ЛУ-П/0.1 предпочтителен с экономической точки зрения по сравнению с материалом на основе углеленты Кулон-500/0.07. Окончательный выбор композиционного материала конструкции (углепластик на основе углеленты ЛУ-П/0.1 и связующего ЭНФБ) был произведён после уточнённого анализа напряжённо"=деформированного состояния с учётом воздействующих факторов на всех этапах эксплуатации и экспериментальных значений характеристик применяемых материалов, определённых при проектировании корпуса.



\Section[n]{Заключение} Следует отметить, что на начальном этапе создания конструкции для определения напряжённо"=деформированного состояния корректно использовать ранее полученные решения аналитических расчётных моделей, что позволяет получить прямую аналитическую связь между искомыми параметрами. Однако роль конечно"=элементной модели размеростабильной конструкции является основной в решении комплексной задачи обеспечения формостабильности и прочности изделия, а современная вычислительная техника и программные продукты позволяют создать подробные модели конструкций и в дальнейшем на основе эксперимента подтвердить их достоверность.



\thanks{{\bf ORCIDs}



{\bf  Владимир Евгеньевич Биткин:}   ??? \url{http://orcid.org/0000-0002-6996-3580}



{\bf Ольга Геннадьевна Жидкова:}  ???  \url{http://orcid.org/0000-0001-9914-8870}



{\bf  Александр Владимирович Денисов:}   ??? \url{http://orcid.org/0000-0002-6996-3580}



{\bf Андрей Викторович Бородавин:}  ???  \url{http://orcid.org/0000-0001-9914-8870}



{\bf Диана Викторовна Митюшкина:}  ???  \url{http://orcid.org/0000-0001-9914-8870}



{\bf  Александр Вениаминович Родионов:}   ??? \url{http://orcid.org/0000-0002-6996-3580}



{\bf Александр Сергеевич Нонин:}  ???  \url{http://orcid.org/0000-0001-9914-8870}



}



\RusRef

\begin{thebibliography}{22}



%1. Углеродные волокна и углекомпозиты: Пер. с англ. / Под ред. Э. Фитцера. -- М.: Мир, 1988. 336 с., ил.

\RBibitem{bit:bib1}

\book Углеродные волокна и углекомпозиты

\ed Э.~Фитцера

\publaddr М.

\publ Мир

\yr 1988

\totalpages 336



%2. Формостабильные и интеллектуальные конструкции из композиционных материалов/ Молодцов Г.А., Биткин В.Е., Симонов В.Ф., Урмансов Ф.Ф. -- М.: Машиностроение, 2000. -- 352 с.

\RBibitem{bit:bib2}

\by Молодцов~Г.~А., Биткин~В.~Е., Симонов~В.~Ф., Урмансов~Ф.~Ф.

\book Формостабильные и интеллектуальные конструкции из композиционных материалов

\publaddr М.

\publ Машиностроение

\yr 2000

\totalpages 352



%3. Климакова Л.А., Половый А.О., Зельнев В.Н. Формостабильность крупногабаритной углепластиковой антенной платформы КА с конструктивно-жесткостной несимметрией. // Механика композиционных материалов и конструкций, 2011, № 2, том. 17. С. 245-254.

\RBibitem{bit:bib3}

\paper Формостабильность крупногабаритной углепластиковой антенной платформы КА с конструктивно"=жесткостной несимметрией

\by Климакова~Л.~А., Половый~А.~О., Зельнев~В.~Н.

\jour Механика композиционных материалов и конструкций

\yr 2011

\issue 2

\vol 17

\pages 245--254



4. Photos: Herschel Space Observatory's Amazing Infrared Images. URL: http://www.space.com/12501-herschel-space-telescope-photos-astronomy.html (дата обращения 28.09.2016).



5. Hubble Team Unveils Most Colorful View of Universe Captured by Space Telescope. URL: http://hubblesite.org/newscenter/archive/releases/2014/27 (дата обращения 18.10.2016).



%6. Безмозгий И.М., Софинский А.Н., Чернягин А.Г. Моделирование в задачах вибропрочности конструкций ракетно-космической техники. // Космическая техника и технологии, 2014, № 3. С. 71-80.

\RBibitem{bit:bib6}

\paper Моделирование в задачах вибропрочности конструкций ракетно"=космической техники

\by Безмозгий~И.~М., Софинский~А.~Н., Чернягин~А.~Г.

\jour Космическая техника и технологии

\yr 2014

\issue 3

\pages 71--80



%7. Биткина Е.В., Денисов А.В., Биткин В.Е. Конструктивно-технологические методы создания размеростабильных космических композитных конструкций интегрального типа. // Механика и машиностроение / Известия Самарского научного центра Российской академии наук, 2012. т. 14, №4(2). -- С. 555-560.

\RBibitem{bit:bib7}

\paper Конструктивно"=технологические методы создания размеростабильных космических композитных конструкций интегрального типа

\by Биткина~Е.~В., Денисов~А.~В., Биткин~В.~Е.

\jour Механика и машиностроение / Известия Самарского научного центра Российской академии наук

\yr 2012

\vol 14

\issue 4(2)

\pages 555--560



%8. Биткин В.Е., Денисов А.В., Денисова M.A., Жидкова О.Г., Назаров Е.В., Рогальская О.И., Мелентьев А.В., Мизинова И.А. Апробирование технологического комплекса изготовления силовых и высокоточных размеростабильных элементов конструкций интегрального типа из волокнистых композиционных материалов. / Известия Самарского научного центра Российской академии наук, 2014. т. 16, №1 (5). -- С. 1320-1327

\RBibitem{bit:bib8}

\paper Апробирование технологического комплекса изготовления силовых и высокоточных размеростабильных элементов конструкций интегрального типа из волокнистых композиционных материалов

\by Биткин~В.~Е., Денисов~А.~В., Денисова~M.~A., Жидкова~О.~Г., Назаров~Е.~В., Рогальская~О.~И., Мелентьев~А.~В., Мизинова~И.~А.

\jour Известия Самарского научного центра Российской академии наук

\yr 2014

\vol 16

\issue 1(5)

\pages 1320--1327



9. MSC.NASTRAN. Quick Reference Guid.



%10. Основы проектирования и изготовления конструкций летательных аппаратов из композиционных материалов: Учебное пособие / В.В. Васильев, A.A. Добряков, A.A. Дудченко и др.- М.: МАИ, 1985. -- 218 с.

\RBibitem{bit:bib10}

\by Васильев~В.~В., Добряков~A.~A. , Дудченко~A.~A.  и др.

\book Основы проектирования и изготовления конструкций летательных аппаратов из композиционных материалов: Учебное пособие

\publaddr М.

\publ МАИ

\yr 1985

\totalpages 218



%11. Васильев В.В. Механика конструкций из композиционных материалов. -- М.: Машиностроение, 1988. -- 272 с.

\RBibitem{bit:bib11}

\by Васильев~В.~В.

\book Механика конструкций из композиционных материалов

\publaddr М.

\publ Машиностроение

\yr 1988

\totalpages 272



%12. Гардымов Г.П., Мешков Е.В., Пчелинцев А.В., Лашманов Г.П., Афанасьев Ю.А. Композиционные материалы в ракетно-космическом аппаратостроении / Под общ. ред. д-ра техн. наук, проф. Г. П. Гардымова и д-ра техн наук, проф. Е. В. Мешкова. - СПб.: СпецЛит, 1999. - 271 с.

\RBibitem{bit:bib12}

\by Гардымов~Г.~П., Мешков~Е.~В., Пчелинцев~А.~В., Лашманов~Г.~П., Афанасьев~Ю.~А.

\book Композиционные материалы в ракетно-космическом аппаратостроении

\eds д-ра техн. наук, проф. Г.~П.~Гардымова и д-ра техн наук, проф. Е.~В.~ Мешкова

\publaddr СПб.

\publ СпецЛит

\yr 1999

\totalpages 271



%13. Комков М.А., Тарасов В.А. Технология намотки композитных конструкций ракет и средств поражения. -- М.: Изд-во МГТУ им. Н.Э. Баумана, 2011. 431 с.

\RBibitem{bit:bib13}

\by Комков~М.~А., Тарасов~В.~А.

\book Технология намотки композитных конструкций ракет и средств поражения

\publaddr М.

\publ Изд-во МГТУ им.~Н.~Э.~Баумана

\yr 2011

\totalpages 431



%14. Воробей В.В., Войтков Н.И. Некоторые прикладные задачи механики  размеростабильных конструкций из композитов. // Механика композитных материалов, 1990, № 2. С. 292-298.

\RBibitem{bit:bib14}

\paper Некоторые прикладные задачи механики  размеростабильных конструкций из композитов

\by Воробей~В.~В., Войтков~Н.~И.

\jour Механика композитных материалов

\yr 1990

\issue 2

\pages 292--298



%15. Зенкевич О. Метод конечных элементов в технике. -- М.: Мир, 1975.

\RBibitem{bit:bib15}

\by Зенкевич~О.

\book Метод конечных элементов в технике

\publaddr М.

\publ Мир

\yr 1975



%16. Якупов Н.М., Киямов Х.Г., Якупов С.Н., Киямов И.Х. Моделирование элементов конструкций сложной геометрии трехмерными конечными элементами. // Механика композиционных материалов и конструкций, 2011, № 1, том. 17. . С. 145-154.

\RBibitem{bit:bib16}

\paper Моделирование элементов конструкций сложной геометрии трехмерными конечными элементами

\by Якупов~Н.~М., Киямов~Х.~Г., Якупов~С.~Н., Киямов~И.~Х.

\jour Механика композиционных материалов и конструкций

\yr 2011

\issue 1

\vol 17

\pages 145--154



%17. Шимкович Д.Г. Расчет конструкций в MSC/NASTRAN for Windows. -- М.: ДМК Пресс, 2003. -- 448 с., ил.

\RBibitem{bit:bib17}

\by Шимкович~Д.~Г.

\book Расчет конструкций в MSC/NASTRAN for Windows

\publaddr М.

\publ ДМК Пресс

\yr 2003

\totalpages 448



%18. Zienkiewicz, O.C.; Taylor, R.L.; Zhu, J.Z. The Finite Element Method: Its Basis and Fundamentals. Elsevier Butterworth-Heinemann, Amsterdam, sixth ed., 2005.

\Bibitem{bit:bib18}

\by Zienkiewicz~O.~C., Taylor~R.~L., Zhu~J.~Z.

\book The Finite Element Method: Its Basis and Fundamentals

\publaddr Amsterdam

\publ Elsevier Butterworth-Heinemann

\yr 2005



%19. Чигарев А.В., Кравчук А.С., Смалюк А.Ф. ANSYS для инженеров // М.: Машиностроение, 2004. 510 с.

\RBibitem{bit:bib19}

\by Чигарев~А.~В., Кравчук~А.~С., Смалюк~А.~Ф.

\book ANSYS для инженеров

\publaddr М.

\publ Машиностроение

\yr 2004

\totalpages 510



%20. Иванов С.Е. Интеллектуальные программные комплексы для технической и технологической подготовки производства /Часть 5. Системы инженерного расчета и анализа деталей и сборочных единиц. Под ред. Куликова Д. Д. Учебно-методическое пособие. -- СПб: СПбГУ ИТМО, 2011. -- 48 с.

\RBibitem{bit:bib20}

\by Иванов~С.~Е.

\book Интеллектуальные программные комплексы для технической и технологической подготовки производства /Часть 5. Системы инженерного расчёта и анализа деталей и сборочных единиц

\ed Куликов~Д.~Д.

\publaddr СПб.

\publ СПбГУ ИТМО

\yr 2011

\totalpages 48



%21. Басов К.А. ANSYS: справочник пользователя. -- М.: ДМК Пресс, 2005 -- 640 с., ил.

\RBibitem{bit:bib21}

\by Басов~К.~А.

\book ANSYS: справочник пользователя

\publaddr М.

\publ ДМК Пресс

\yr 2005

\totalpages 640



22. Оптико-электронные системы для дистанционного зондирования Земли. URL: http://www.lomo-tech.ru/photos/lomo{\_}kosm{\_}otkr.pdf (дата обращения 20.09.2016)



\end{thebibliography}

\RuDate









\newpage



\makeentitle



\thanks{{\bf Declaration of Financial and Other Relationships.}

This work was supported by the program of the Ural Branch of RAS (project no.~15--10--1--18).

Each author has participated in the article concept development and in the manuscript writing.

The authors are absolutely responsible for submitting the final manuscript in print.

Each author has approved the final version of manuscript.

The authors have not received any fee for the article.

}



\pagebreak





\thanks{{\bf ORCIDs}



{\bf Vladimir E. Bitkin:}  ???  \url{http://orcid.org/0000-0002-6996-3580}



{\bf Olga G. Zhidkova:}  ???  \url{http://orcid.org/0000-0001-9914-8870}



{\bf Alexandr V. Denisov:}  ???  \url{http://orcid.org/0000-0002-6996-3580}



{\bf Andrey V. Borodavin:}  ???  \url{http://orcid.org/0000-0001-9914-8870}



{\bf Diana V. Mityushkina:}  ???  \url{http://orcid.org/0000-0001-9914-8870}



{\bf Alexandr V. Rodionov:}  ???  \url{http://orcid.org/0000-0002-6996-3580}



{\bf Alexandr S. Nonin:}  ???  \url{http://orcid.org/0000-0001-9914-8870}



}





\EngRef

\begin{thebibliography}{40}



\Bibitem{mult:Ref1:en}

\paper Damping of structural vibrations with piezoelectric

 materials and passive electrical networks

\by Hagood~N.\,W., von Flotow~A.

\jour J. Sound Vib.

\yr 1991

\vol 146

\pages 243--268

\crossref{10.1016/0022-460X(91)90762-9}



\Bibitem{mult:Ref2:en}

\by Moheimani~S.\,O.\,R., Fleming~A.\,J.

\book Piezoelectric transducers for vibration control and damping

\publ Springer Science and Business Media

\yr 2006

\totalpages 272



\Bibitem{mult:Ref3:en}

\paper Piezoelectric shunts with parallel R–L circuit for

structural damping and vibration control

\by Wu~S.\,Y.

\jour Proc. SPIE  2720, Smart Structures and Materials 1996:

Passive Damping and Isolation

\yr 1996

\pages 259--269

\crossref{10.1117/12.239093}



\Bibitem{mult:Ref4:en}

\paper Multimodal passive vibration suppression with

piezoelectric materials and resonant shunts

\by Hollkamp~J.\,J.

\jour J. Intell. Mater. Syst. Struct.

\yr 1996

\issue 5

\pages 49--56

\crossref{10.1177/1045389X9400500106}



\Bibitem{mult:Ref5:en}

\paper Method for multiple mode shunt damping of structural

vibration using a single PZT transducer

\by Wu~S.\,Y.

\jour Proc. SPIE  3327, Smart Structures and Materials 1998:

Passive Damping and Isolation

\yr 1998

\pages 159--168

\crossref{10.1117/12.310680}



\Bibitem{mult:Ref6:en}

\paper Multiple PZT transducers implemented with multiple-mode

piezoelectric shunting for passive vibration damping

\by Wu~S.\,Y.

\jour Proc. SPIE  3672, Smart Structures and Materials 1999:

Passive Damping and Isolation

\yr 1999

\pages 112--122

\crossref{10.1117/12.349774}



\Bibitem{mult:Ref7:en}

\paper Structural vibration damping experiments using

improved piezoelectric shunts

\by Wu~S.\,Y., Bicos~A.\,S.

\jour Proc. SPIE  3045, Smart Structures and Materials 1997:

Passive Damping and Isolation

\yr 1997

\pages 40--50

\crossref{10.1117/12.274217}



\Bibitem{mult:Ref8:en}

\paper Current flowing multiple mode piezoelectric shunt dampener

\by Behrens~S.\, Moheimani~S.\,O.\,R.

\jour Proc. SPIE 4697, Smart Structures and Materials 2002:

Damping and Isolation

\yr 2002

\pages 117--127

\crossref{10.1117/12.472658}



\Bibitem{mult:Ref9:en}

\paper Multiple mode passive piezoelectric shunt dampener

\by Behrens~S., Moheimani~S.\,O.\,R.\, Fleming~A.\,J.

\jour Proc. IFAC Mechatronics

\yr 2002

\pages



\Bibitem{mult:Ref10:en}

\paper Adaptive piezoelectric shunt damping

\by Fleming~A.\,J., Moheimani~S.\,O.\,R.

\jour Smart Mater. Struct.

\yr 2003

\vol 12

\issue 1

\pages 36—-48



\Bibitem{mult:Ref11:en}

\paper Optimal resistive elements for multiple mode shunt damping of a piezoelectric laminate beam

\by Behrens~S., Moheimani~S.\,O.\,R.

\jour Proceedings of the 39th IEEE Conference on Decision and Control

\yr 2000

\vol 4

\pages 4018--4023

\crossref{10.1109/CDC.2000.912343}



\Bibitem{mult:Ref12:en}

\paper A current-flowing electromagnetic shunt damper for multi-mode

vibration control of cantilever beams

\by Cheng~T.\,H., Oh~I.\,K.\,A.

\jour Smart Materials and Structures

\yr 2009

\vol 18

\issue 9

%\finalinfo Article ID 095036

\crossref{10.1088/0964-1726/18/9/095036}



\Bibitem{mult:Ref13:en}

\paper Vibration control in plates by uniformly distributed

PZT actuators interconnected via electric networks

\by Vidoli~S., dell'Isola~F.

\jour European Journal of Mechanics - /A Solids

\yr 2001

\vol 20

\issue 3

\pages 435--456

\crossref{10.1016/S0997-7538(01)01144-5}



\Bibitem{mult:Ref14:en}

\paper Circuit analog of a beam and its application to

multimodal vibration damping, using piezoelectric transducers

\by Porfiri~M., dell'Isola~F., Mascioli~F.\,M.

\jour Int. J. Circ. Theor. Appl.

\yr 2004

\vol 32

\issue 4

\pages 167--198

\crossref{10.1002/cta.273}



\Bibitem{mult:Ref15:en}

\paper Piezoelectromechanical structures: new trends towards

the multimodal passive vibration control

\by dell'Isola~F., Henneke~E.\,G., Porfiri~M.

\jour Proceedings of SPIE Vol. 5052, Smart Structures

and Materials 2003: Damping and Isolation

\yr 2003

\pages 392--402

\crossref{10.1117/12.483803}



\Bibitem{mult:Ref16:en}

\paper Comparison of piezoelectronic networks acting as

distributed vibration absorbers

\by Maurini~C., dell'Isola~F., Del Vescovo~D.

\jour Mechanical Systems and Signal Processing

\yr 2004

\vol 18

\issue 5

\pages 1243--1271

\crossref{10.1016/S0888-3270(03)00082-7}



\Bibitem{mult:Ref17:en}

\paper Multimode vibration control using several piezoelectric

transducers shunted with a multiterminal network

\by Giorgio~I., Culla~A., Del Vescovo~D.

\jour Arch. Appl. Mech

\yr 2009

\vol 79

\issue 9

\pages 859--879

\crossref{10.1007/s00419-008-0258-x}



\Bibitem{mult:Ref18:en}

\paper Multimodal Vibration Damping through Piezoelectric

Patches and Optimal Resonant Shunt Circuits

\by Viana~F.\,A.\,C., Steffen~V.\,Jr.

\jour J. of the Braz. Soc. of Mech. Sci. and Eng

\yr 2006

\vol XXVIII

\issue 3

\pages 293--310

\crossref{10.1590/S1678-58782006000300007}



\Bibitem{mult:Ref19:en}

\paper Broadband vibration control through periodic arrays

of resonant shunts: experimental investigation on plates

\by Casadei~F., Ruzenne~M., Dozio~L., Cunefare~K.\,A.

\jour Smart Materials and Structures

\yr 2010

\vol 19

\issue 1

%\finalinfo Article ID 015002

\crossref{10.1088/0964-1726/19/1/015002}



\Bibitem{mult:Ref20:en}

\by Washizu~K.

\book Variational Methods in Elasticity and Plasticity

\publaddr London

\publ Pergamon Pr.

\yr 1982

\totalpages 540



\Bibitem{mult:Ref21:en}

\by Parton~V.\,Z., Kudryavtsev~B.\,A.

\book Electromagnetoelasticity: Piezoelectrics and

Electrically Conductive Solids

\publaddr New York

\publ Gordon and Breach Science Publishers Ltd

\yr 1988

\totalpages 503



\Bibitem{mult:Ref22:en}

\by Karnaukhov~V.\,G., Kirichok~I.\,F.

\book Electrothermal viscoelasticity

\publaddr Kiev

\publ Nauk. dumka

\yr 1988

\totalpages 319

%\finalinfo (In Russian)



\Bibitem{mult:Ref23:en}

\paper Simulation and optimization of dynamic characteristics

of piezoelectric smart structures

\by Matveenko~V.\,P., Kligman~E.\,P., Yurlov~M.\,A., Yurlova~N.\,A.

\jour Phys. Mesomech.

\yr 2012

\vol 15

\issue 3-4

\pages 190--199

\crossref{10.1134/S1029959912020063}



\Bibitem{mult:Ref24:en}

\paper Optimization of the damping properties electro-viscoelastic

objects with external electric circuits

\by Matveenko~V.\,P., Yurlov~M.\,A., Yurlova~N.\,A.

\inbook Mechanics of Advanced Materials: Analysis of

Properties and Performance

\eds Silberschmidt~V.\,V., Matveenko~V.\,P.

\publaddr Switzerland

\publ Springer International Publishing

\yr 2015

\pages 79--100

\crossref{10.1007/978-3-319-17118-0-4}



\Bibitem{mult:Ref25:en}

\paper The optimal placement of the piezoelectric element in a

structure based on the solution of the problem of natural vibrations

\by Sevodina~N.\,V., Yurlova~N.\,A., Oshmarin~D.\,A.

\jour Solid State Phenomena

\yr 2016

\vol 243

\pages 67--74

\crossref{10.4028/www.scientific.net/SSP.243.67}



\Bibitem{mult:Ref26:en}

\paper Justification of equivalent substitution circuits

used to optimize the dissipative properties of electroelastic

bodies with external electric circuits

\by Ivanov~A.\,S. Matvienko~V.\,P. Oshmarin~D.\,A, Sevodina~N.\,V., Yurlov~M.\,A., Yurlova~N.\,A.

\jour Mechanics of Solids

\yr 2016

\vol 51

\issue 3

\pages 273--283

\crossref{:10.3103/S0025654416030043}



\end{thebibliography}



\EngDate





\end{document}



 















